% !TEX TS-program = pdflatexmk
\documentclass[11pt]{amsart}

%\usepackage{geometry}                % See geometry.pdf to learn the layout options. There are lots.
%\geometry{letterpaper}                   % ... or a4paper or a5paper or ... 
%\geometry{landscape}                % Activate for for rotated page geometry
%\usepackage[parfill]{parskip}    % Activate to begin paragraphs with an empty line rather than an indent

\usepackage[margin=1in]{geometry}

\usepackage{amsmath,amssymb,amsthm, latexsym, mathrsfs, pdfsync, multicol,
%setspace,
%graphics, 
fancybox, fancyhdr,
graphicx, enumerate,
subfig, tikz, pgfplots,array}

\usetikzlibrary{calc}
\usetikzlibrary{shapes}

\pgfplotsset{my style/.append style={axis x line=middle, axis y line=middle, 
xlabel=$x$,ylabel=$y$,
every axis x label/.style={
    at={(ticklabel* cs:1)},
    anchor=west,
},
every axis y label/.style={
    at={(ticklabel* cs:1)},
    anchor=south,
},}}



\usepackage[dvipsnames,usenames]{xcolor}
\usepackage[mathscr]{euscript}
\usepackage[%final, 
colorlinks=true, backref=page%,
%linkbordercolor=red, pdfborderstyle={/S/U/W 1}
]{hyperref}
\usepackage[normalem]{ulem}
\usepackage{lineno}
\DeclareGraphicsRule{.tif}{png}{.png}{`convert #1 `dirname #1`/`basename #1 .tif`.png}

\usepackage{cleveref}
\newtheorem{theorem}{Theorem}
\newtheorem{corollary}[theorem]{Corollary}
\newtheorem{conjecture}[theorem]{Conjecture}
\newtheorem{lemma}[theorem]{Lemma}
\newtheorem{proposition}[theorem]{Proposition}
\Crefname{definition}{Definition}{Definitions}

\theoremstyle{definition}
\newtheorem{definition}[theorem]{Definition}
\newtheorem{axiom}{Axiom}

\Crefname{exercise}{Exercise}{Exercises}
\newtheorem{exercise}{Exercise}%[section]
\theoremstyle{remark}
\newtheorem{remark}{Remark}
\newtheorem{example}[theorem]{Example}
\newcommand\RR{\ensuremath{\mathbb R}}
\newcommand\NN{\ensuremath{\mathbb N}}
\newcommand\ZZ{\ensuremath{\mathbb Z}}
\newcommand\QQ{\ensuremath{\mathbb Q}}
\newcommand\CC{\ensuremath{\mathbb C}}
\newcommand\EE{\ensuremath{\mathbb E}}
\newcommand\bc{\begin{center}}
\newcommand\ec{\end{center}}
\newcommand\be{\begin{enumerate}}
\newcommand\ee{\end{enumerate}}
\newcommand\bi{\begin{itemize}}
\newcommand\ei{\end{itemize}}
\newcommand\bs{\begin{slide}}
\newcommand\es{\end{slide}}
\newcommand\bx{\begin{exercise}}
\newcommand\ex{\end{exercise}}
\newcommand{\ol}[1]{\overline{#1}}
\newcommand{\oimp}[1]{\overset{#1}{\Longleftrightarrow}}
\newcommand{\bv}[1]{\ensuremath{ \vec{\mathbf{#1}}} }
\newcommand{\mc}[1]{\ensuremath{\mathcal{#1}}}
\newcommand{\ARP}{\cite{McMSch02}}
%\DeclareMathOperator{\Aut}{Aut} 
\newcommand{\Aut}[1]{\ensuremath{{\mathscr Aut}(#1)}}
\newcommand{\Con}[1]{\ensuremath{{\mathscr Con}(#1)}}
\newcommand{\Mod}[1]{\ (\text{mod}\ #1)}
\newcommand{\Mon}[1]{\ensuremath{{\mathscr Mon}(#1)}}
\newcommand{\Pyr}[1]{\ensuremath \rm{ Pyr} (#1)}
\newcommand{\mix}{\ensuremath{\diamondsuit}}
\newcommand{\covers}{\ensuremath{\searrow}}
\newcommand{\normale}{\vartrianglelefteq}
\newcommand{\normal}{\vartriangleleft}
\newcommand{\nin}{\not\in}

\renewcommand{\d}{\displaystyle}

\theoremstyle{comment}
\newtheorem*{acomment}{\color{BrickRed}{Comment}}
\newcommand{\acom}[2]{{\begin{acomment}[#1]\color{OliveGreen} #2 \color{black} \end{acomment}}}
\newcommand{\red}[1]{{\color{BrickRed} #1}}
\newcommand{\blue}[1]{{\color{MidnightBlue}#1}}
\newcommand{\green}[1]{{\color{OliveGreen} #1 }}
\newcommand{\mwrong}[1]{\red{\cancel{#1}}}
\newcommand{\wrong}[2]{\red{\sout{#1} {#2}}}
\definecolor{OliveGreen}{rgb}{.3,.5,.2}
\definecolor{MidnightBlue}{rgb}{.3,.4,.6}

\title{MT 1 Review Exercises}

%\author{Gordon Williams}
%\author{Course}
\begin{document}
\bc{\Large MT 1 Review Exercises}\ec
This is a selection of exercises drawn from previous exams for you to practice on.
\begin{enumerate}
\item (10 points) The graphs of two functions $f(x)$, shown thick, and $g(x)$, shown dashed, are given below.  %whose domain is $(-6, 7
Determine the following. If the value does not exist, write ``DNE''.
%
%\begin{figure}[ht]
\begin{center}
\begin{tikzpicture}[scale=.9]
\draw[ultra thin] (-6.9,-3) grid (6.9,3.1);
\draw[thick,<->] (-7.2,0) -- (7.2,0);
\draw[thick,<->] (0,-3.1) -- (0,3.1);
%\draw[dashed] (7,-2) -- (7,3);
%\draw[dashed] (-6,-2)--(-6,3);

%%%%% f(x) %%%%
\draw [ultra thick] (-5, -1) parabola bend (-4,2)  (-3,0) to  (-2,-1) to (1,-1);
\draw [ultra thick, ->] (1, -2)  parabola bend (2,-2) (3,-1)-- (4,1) -- (6,3);
\draw[fill=white, thick] (4,1) circle  (1.2 mm);
%\draw[fill=white] (6, 2) circle (1.2mm);
\draw[fill=black] (1,-2) circle (1.2mm);
\draw[fill=white] (1,-1) circle (1.2mm);
\draw[fill=white] (-5,-1) circle (1.2mm);
%\draw[fill=black] (6,3) circle (1.2mm);

\draw (2.9, -1.5) node[right] {$f(x)$};
\draw (3,2) node[above left] {$g(x)$};

%%%%% g(x) %%%%%
%\draw[ultra thick, dashed, <->](-6.8, -2.9)-- (-6, -2.9)  to[out = 5, in = 180+45](-3,-2) to[bend left = 30] (-1,0) to [out = 45, in = -90] (1.9,2.5)--(1.9,3);

\draw[ultra thick, dashed, <->] plot [smooth] coordinates { (-6.8, -2.9) (-5, -2.8) (-4, -2.6) (-3,-2) (0, .2) %(.5,0) 
(2.4,1) 
(2.95,3)};

\foreach \i in {-6, -5, ..., 6}
{\draw (\i,0) node[below] {\i};
%\draw (\i+.5, -.1) -- (\i+.5, .1);
}
\foreach \i in { -3, -2, -1,1,2, 3}
{\draw (0,\i) node[left] {\i};}
%\foreach \i in { -2, -1,0,1,2}{
%\draw (-.1, \i+.5) -- (.1,\i+.5);}
\end{tikzpicture}

\end{center}
%\end{figure}

\newcommand{\ans}{\rule{1.5cm}{.5 pt}}


\begin{enumerate}
\begin{multicols}{2}

\item $\d\lim_{x\to -5^{+}} f(x) = \ans$


\item $ \d \lim_{x\to -3} f(x) = \ans$


\item $\d \lim_{x \to \infty} f(x) = \ans$



\item $\d \lim_{x\to 1} f(x) = \ans $


\item $f'(-4) = \ans $


\item $\d\lim_{x\to 3^{-}} g(x) = \ans$


\item $\d \lim_{x \to -\infty} g(x) = \ans$






\item $\d \lim_{x\to -3} f(g(x)) = \ans $



\end{multicols}

\item What is the domain of $f(x)$? Give your answer in interval form.\\



\item Where in the domain of $f(x)$ is the function continuous? Give your answer in interval form.



\end{enumerate}



\item Evaluate the following limits. Justify your
answers with words and/or any relevant algebra. Be sure to use proper
notation, as points will be deducted for not doing so. 


\be
\item  $\d \lim_{x \to -5} \frac{x^2-25}{x^2+2x-15}$

\item $\d \lim_{t \to -3} \frac{6+4t}{t^2+1}$
% $\d \lim_{x \to 5^-} \frac{-\sqrt x}{(x-4)^3}$
% $\d \lim_{x \to 4^+} \frac{-\sqrt x}{(4-x)^3}$

% $\d \lim_{x \to 3} \frac{ \frac 1{x^2} - \frac 1 4}{x-2}$
%  $\d \lim_{x \to 3} \frac{ \frac 1{x^2} - \frac 1 {16}}{x-4}$

\item This is a two parter:
\be[(i)] 
\item $ \d \lim_{x \to -1^-} \sqrt{x^2-1}$
\item Why do we not evaluate $ \d \lim_{x \to -1^+} \sqrt{x^2-1}$ ? Explain using a sentence.
\ee
%% $ \d \lim_{x \to - \infty} \frac{\sqrt{1 + 9x^8}}{5+ x^4}$
%% $ \d \lim_{x \to - \infty} \frac{\sqrt{7 + 25x^6}}{2+x^3}$
%%
%\ee
%
\ee

\newpage
\item Use the graph of $f(x)$ (on top) to sketch the graph of $f'(x)$ (on bottom).\\

\begin{tikzpicture}
\begin{axis}[xscale=3/2,yscale=1.2, thick, my style, xtick={-4,-3,...,8}, ytick={-4,-3,...,5},
xmin=-4.2, xmax=8.5, ymin=-4.2, ymax=5.5, minor y tick num=1,
        minor x tick num=1, mark size=3.0pt]
        \addplot[domain=-3:-1, smooth, ultra thick,<-] coordinates { (-3,4) (-1,0)};
        \addplot[domain=-1:2, smooth, ultra thick] coordinates {  (-1,0) (2,3)};
        \addplot[domain=2:8, ultra thick,->] {3/(x-1)};
        \foreach \i in {-4,-3,...,8}{
	\addplot[dotted, gray, domain=-4:8]{\i};}
	\foreach \i in {-4,-3,...,8}{
	\addplot[dotted, gray] coordinates {(\i,-4) (\i,5)};
	}
\end{axis}
\node at (0,5){{\Large{$f(x)$}}};
\end{tikzpicture}
    
   
\begin{tikzpicture}
\begin{axis}[xscale=3/2,yscale=1.2, thick, my style, xtick={-4,-3,...,8}, ytick={-4,-3,...,5},
xmin=-4.2, xmax=8.5, ymin=-4.2, ymax=5.5, minor y tick num=1,
        minor x tick num=1, mark size=3.0pt]
        \foreach \i in {-4,-3,...,8}{
	\addplot[dotted, gray, domain=-4:8]{\i};}
	\foreach \i in {-4,-3,...,8}{
	\addplot[dotted, gray] coordinates {(\i,-4) (\i,5)};
	}
\end{axis}
\node at (0,5){{\Large{$f'(x)$}}};
\end{tikzpicture}


\item  An icicle is growing
at the back of a house.  At time $t\ge 0$
days the length of the icicle is
$$
\ell(t) = \frac{14 t + 4}{t+2}
$$
inches.  Answer the following questions about the
function $\ell(t)$ and be sure to include \textbf{units}
in your answers.
\be
\item Compute the average rate of change of
the length of the icicle from time $t=0$ to time $t=2$
days.  \textbf{Do not forget units!}


\item  It is a known fact that $\d \ell'(t) = \frac{24}{(t+2)^2}$.
Compute $\ell'(2)$ (\textbf{with units}) and explain what this quantity 
means about the icicle.


\item 
Compute $\lim_{t\to\infty} \ell(t)$.
Then explain what this quantity means in precise
but everyday language that the general public would understand. 
Do not forget units!
\ee

\newpage
\item Is the following function continuous at $x=3$? Justify your answer. 

$
f(x) = \begin{cases}
\frac{6-2x}{x-3} & \text{if } x \leq 3\\
2  & \text{if } x =3\\
x-5  & \text{if } x \geq 3
\end{cases}
$

\item 
Let $H(t)$ be the daily cost, in dollars, of heating a building when the outside temperature is $t$ degrees Fahrenheit.

	\begin{enumerate}

	\item What is the meaning of $H(0)=50$? Your answer should be a complete sentence and must include units.



	\item What are the units of $H'(t)$?


	\item Do you expect $H'(t)$ to be positive or negative? \textbf{Justify} your answer using complete sentences. 

	\end{enumerate}
\item 
 Consider the function \[\displaystyle f(x)=\frac{1}{6-x}.\] 

Find $f'(5)$ using the {\bf limit definition of the derivative} and show your work using all appropriate notation. 

No credit will be awarded for using other methods. Begin by writing down the limit definition of the derivative.

\item  For each of the following functions, compute the derivative. 

{\bf You do not need to simplify your answers.}
 
\be

\item $y = 11x^{3} -\frac{4x}{3} + x^{2} + x^{5/3}$

 
\item $\displaystyle{ a(\theta) = \theta ^3 \cos (\theta)}$ %% product rule
  

  
\item $\displaystyle{ f(t) = t \sqrt{t} - \frac{1}{8t^4} + \sqrt{2} }$ %% quotient rule or cleverness
  
  

%\item $\displaystyle{ y = \frac{\sin x}{\sin x - \cos x} }$ This is a great problem but we already tested -- twice -- if they know the derivative of sine

%\item $\ds{ y = \frac{ 3 x^{2} - 7x^{3}}{4x^{5} - 6 }}$ %% obligatory quotient rule
%
%


\ee

%\item A company's profit in thousands of dollars when $x$ units are produced is given by $$P(x)=-0.001x^2+0.22x-42.$$
%
%	\begin{enumerate}
%
%	\item Observe that $P(1200) = 78$. Interpret this fact in the context of the problem. To earn full credit your answer should be a complete sentence and must include units.
%
%	
%
%	\item Observe that $P'(1200) = -0.02$. Interpret this fact in the context of the problem. To earn full credit your answer should be a complete sentence and must include units.
%
%	
%
%	\item Suppose the company is currently producing 1200 units. Should the company increase production if it wants to increase profits? \textbf{Justify} your answer.
%
%	\end{enumerate}

\item The questions below reference the function $f(x)=(3-x)(x-7).$ 
	\begin{enumerate}
	\item Find the slope of the line tangent to at $x=4$. 
	
	
	\item Sketch the tangent line to $f(x)$ at $x=4$ on the graph of $f(x)$ below.

\begin{center}
 \begin{tikzpicture}
%quadratic
\begin{axis}[scale=1,xscale=1, thick, my style, xtick={0,1,...,10}, ytick={-6,-5,...,5},xmin=-0.5, xmax=9.5, ymin=-6.5, ymax=5.5, 
mark size=3.0pt, grid = major]
\addplot[ultra thick, <->,domain=0:9, samples=100]{(3-x)*(x-7)};
%\addplot[ultra thick, <-,domain=-6:0, samples=100]{-x};
\node at (9,2){$f(x)$};
\end{axis}
\end{tikzpicture}\end{center}
\end{enumerate}


\end{enumerate}

 \end{document}
 \end

  